%=====================================================================
\chapter{Fuzzy Logic}\label{ch:fuzzy}
%=====================================================================
\locator{5}{Part V --- Reasoning Under Uncertainty}

\begin{outcomes}
By the end of this chapter you will be able to:
\begin{tight}
  \item \textbf{Distinguish} vagueness from uncertainty, and \textbf{say} which of
        the two fuzzy logic addresses.
  \item \textbf{Define} membership functions and \textbf{apply} the fuzzy
        operators.
  \item \textbf{Run} a Mamdani inference by hand: fuzzify, evaluate the rules,
        clip, aggregate, defuzzify.
  \item \textbf{Explain} why a fuzzy controller's output is smooth where a
        threshold rule's is not, and why that matters.
  \item \textbf{State} the limitations of centroid defuzzification and of the
        method generally.
\end{tight}
\end{outcomes}

\begin{prereq}
\chref{uncertainty} --- and specifically the idea of a degree of belief, which is
what a degree of \emph{membership} is repeatedly confused with. This chapter branches
from Chapter~14 rather than following \chref{bayesnets}: it answers a different
question.
\end{prereq}

\begin{keyterms}
vagueness $\bullet$ fuzzy set $\bullet$ membership function $\bullet$ degree of
membership $\bullet$ linguistic variable $\bullet$ fuzzy partition $\bullet$ fuzzy
operators $\bullet$ fuzzification $\bullet$ firing strength $\bullet$ Mamdani
inference $\bullet$ clipping $\bullet$ aggregation $\bullet$ defuzzification
$\bullet$ centroid
\end{keyterms}

\section{Vagueness Is Not Uncertainty}

Two sentences that sound alike and are not:

\begin{quote}
\emph{This project is probably high-risk.}\\
\emph{This project is fairly high-risk.}
\end{quote}

The first is \term{uncertainty}. There is a fact of the matter --- the project either
is or is not high-risk --- and the speaker does not know which.
\chref{uncertainty} is about that, and evidence resolves it: learn enough and the
probability goes to $0$ or $1$.

The second is \term{vagueness}. The speaker knows everything there is to know about
the project. The trouble is that ``high-risk'' has no sharp boundary, and no amount
of evidence will produce one. A project at $70$ effort-days with moderate
uncertainty is high-risk to a degree, and that is the whole truth about it.

\begin{definitionbox}[Fuzzy set and membership]
A \term{fuzzy set} $A$ over a domain is given by a \term{membership function}
$\mu_{A}(x) \in [0,1]$ saying to what degree $x$ belongs to $A$.

$\mu_{\mathit{high}}(70) = 0.4$ does \emph{not} mean ``a $40\%$ chance the project is
high-risk''. It means the project is high-risk to degree $0.4$ --- a statement about
the \emph{word}, not about the world, and it does not change when more evidence
arrives.
\end{definitionbox}

\begin{alertbox}[The confusion this chapter exists to prevent]
A probability and a membership are both numbers in $[0,1]$ and they mean entirely
different things.

$P(\mathit{defect}) = 0.3$ describes ignorance about a sharp fact; test the release
and it resolves to $0$ or $1$. $\mu_{\mathit{high}}(70) = 0.4$ describes a sharp
fact about a blurry word; measure the project more carefully and it stays $0.4$.

Probabilities over a partition sum to one and memberships do not. Combining
probabilities uses multiplication and Bayes' rule; combining memberships uses
$\min$ and $\max$. And Russell and Norvig are openly sceptical that the second
apparatus is ever necessary --- their position is that a well-specified
probabilistic model can do this work. It is worth reading their argument, and
worth knowing that fuzzy controllers run production systems from cement kilns to
camera autofocus, which is an argument of a different kind.
\end{alertbox}

\section{Membership Functions and Linguistic Variables}

A \term{linguistic variable} takes values that are words --- \emph{low}, \emph{medium},
\emph{high} --- each defined by a membership function over the underlying numeric
domain.

The usual construction is a \term{fuzzy partition}: overlapping triangles whose peaks
sit at evenly spaced points and whose memberships sum to one everywhere. Two
consequences: every input is described by at most two words, and the description
changes continuously as the input moves.

\dsfig{%
\begin{tikzpicture}
\begin{axis}[
  width=12.4cm, height=5.0cm,
  xlabel={requirement uncertainty (team's 0--10 scale)},
  ylabel={degree of membership $\mu$},
  xlabel style={font=\scriptsize}, ylabel style={font=\scriptsize},
  tick label style={font=\tiny},
  xmin=0, xmax=10, ymin=0, ymax=1.15,
  grid=both, major grid style={gray!18}, minor grid style={gray!8},
  legend style={font=\tiny, at={(0.5,1.24)}, anchor=north, legend columns=3,
                draw=gray!40, fill=white},
  legend cell align=left]

\addplot[primary, line width=1.3pt] coordinates {(0,1) (5,0) (10,0)};
\addlegendentry{low}
\addplot[goldhl, line width=1.3pt] coordinates {(0,0) (5,1) (10,0)};
\addlegendentry{medium}
\addplot[accent, line width=1.3pt] coordinates {(0,0) (5,0) (10,1)};
\addlegendentry{high}

\draw[greenhl, line width=1pt, dash pattern=on 3pt off 2pt]
  (axis cs:4,0) -- (axis cs:4,1.05);
\addplot[only marks, mark=*, mark size=1.7pt, primary]
  coordinates {(4,0.2)};
\addplot[only marks, mark=*, mark size=1.7pt, goldhl]
  coordinates {(4,0.8)};
\node[font=\tiny, text=greenhl!50!black, anchor=west, align=left]
  at (axis cs:4.2,1.0) {uncertainty 4 is};
\node[font=\tiny, text=greenhl!50!black, anchor=west, align=left]
  at (axis cs:4.2,0.86) {\emph{low} to degree 0.2,};
\node[font=\tiny, text=greenhl!50!black, anchor=west, align=left]
  at (axis cs:4.2,0.72) {\emph{medium} to degree 0.8};
\end{axis}
\end{tikzpicture}}%
{A fuzzy partition of the uncertainty scale. The memberships sum to one at every
point, so each input is described by at most two words and the description moves
smoothly as the input does.}{memberships}

\section{The Fuzzy Operators}

Conjunction, disjunction and negation extend to memberships in the standard way:
\[
  \mu_{A \cap B}(x) = \min\bigl(\mu_{A}(x), \mu_{B}(x)\bigr), \quad
  \mu_{A \cup B}(x) = \max\bigl(\mu_{A}(x), \mu_{B}(x)\bigr), \quad
  \mu_{\neg A}(x) = 1 - \mu_{A}(x)
\]

Note what is \emph{not} true. In classical logic $A \cap \neg A = \varnothing$; here
$\min(\mu, 1-\mu)$ is $0.5$ when $\mu = 0.5$, so a borderline project is in
``high-risk and not high-risk'' to degree $0.5$. That is a feature of the intended
reading --- a genuinely borderline case is half in and half out of a vague category
--- and it is also why fuzzy logic is not a logic in the sense of
\chref{propositional}: there is no entailment relation and no soundness theorem.

\section{Mamdani Inference}

The classical fuzzy controller, in five steps. The example is the team's risk model:
inputs are \emph{effort} (person-days) and \emph{requirement uncertainty}, the
output is \emph{risk}, and there are nine rules --- one for each combination of two
linguistic variables with three values each.

\dsfig{%
\begin{tikzpicture}[font=\scriptsize,
  st/.style={rounded corners=3pt, draw=#1, line width=1.1pt, fill=#1!8,
             text=#1!50!black, text width=2.45cm, align=flush left,
             inner sep=5pt, minimum height=2.4cm},
  hd/.style={rounded corners=2pt, fill=#1, text=white,
             font=\tiny\bfseries, inner sep=2.5pt, text width=2.3cm,
             align=center},
  a/.style={-{Stealth[length=2mm]}, neutral!70, line width=1pt}]

\node[st=primary]   (s1) at (0,0)     {\\[2pt]effort $=45$\\ uncertainty $=4$\\[4pt]
  become\\[2pt] \emph{low} 0.1,\\ \emph{medium} 0.9\\[2pt]
  \emph{low} 0.2,\\ \emph{medium} 0.8};
\node[st=secondary] (s2) at (3.1,0)   {\\[2pt]each rule's strength is the
  $\min$ of its conditions\\[4pt]
  4 of the 9 rules have non-zero strength};
\node[st=goldhl]    (s3) at (6.2,0)   {\\[2pt]each rule \emph{clips} its
  output word at its own strength\\[4pt]
  \emph{low} at 0.2\\ \emph{medium} at 0.8};
\node[st=greenhl]   (s4) at (9.3,0)   {\\[2pt]aggregate the clipped shapes
  by $\max$\\[4pt] one combined shape over the risk scale};
\node[st=accent]    (s5) at (12.4,0)  {\\[2pt]take the centroid of that
  shape\\[4pt] \textbf{risk $=49.04$}};

\node[hd=primary]   at (s1.north) {1. FUZZIFY};
\node[hd=secondary] at (s2.north) {2. FIRE RULES};
\node[hd=goldhl]    at (s3.north) {3. CLIP};
\node[hd=greenhl]   at (s4.north) {4. AGGREGATE};
\node[hd=accent]    at (s5.north) {5. DEFUZZIFY};

\foreach \x/\y in {s1/s2,s2/s3,s3/s4,s4/s5}{\draw[a] (\x) -- (\y);}

\node[rounded corners=3pt, fill=lightbg, draw=primary!30, line width=0.8pt,
      text=primary!70!black, font=\scriptsize, text width=13.6cm,
      align=flush left, inner sep=6pt] at (6.2,-2.1)
  {\textbf{Steps 1 and 5 are the interface; steps 2 to 4 are the reasoning.} The
   inputs and the output are ordinary numbers, so the controller drops into an
   existing system unchanged --- and in between, the knowledge is nine sentences a
   delivery lead wrote in English. That combination is why fuzzy control was
   adopted in industry far more widely than in research.};
\end{tikzpicture}}%
{Mamdani inference on the project \emph{Beacon}, effort $45$ and uncertainty $4$.
Four of the nine rules fire, and the output is a single number.}{mamdani}

\begin{examplebox}[The nine rules]
\begin{tabular}{@{}L{4.5cm}L{4.5cm}L{5.4cm}@{}}
\toprule
\textbf{Effort} & \textbf{Uncertainty} & \textbf{Risk} \\
\midrule
low & low & low \\
low & medium & low \\
low & high & medium \\
medium & low & low \\
medium & medium & medium \\
medium & high & high \\
high & low & medium \\
high & medium & high \\
high & high & high \\
\bottomrule
\end{tabular}

\smallskip
Nine rules for two inputs with three values each. Note how this scales: three inputs
would need $27$, four would need $81$. \emph{Rule explosion} is the practical limit
on fuzzy control, and it is why these systems typically have two or three inputs.
\end{examplebox}

\section{Defuzzification}

Aggregation leaves a shape; a controller needs a number. The \term{centroid} method
takes the shape's centre of area:
\[
  x^{*} \;=\; \frac{\int x\, \mu(x)\, dx}{\int \mu(x)\, dx}
\]
computed in practice by summing over a grid.

\begin{alertbox}[Centroid defuzzification cannot reach the extremes]
Run the risk model across its whole input range and the output never leaves roughly
$[16.6, 83.4]$ --- it never says $0$ and never says $100$, even for a project with
minimum effort and minimum uncertainty.

The reason is geometric rather than conceptual: the centroid of a shape clipped from
a triangle centred at $0$ still has area to the right of $0$, so its centre of area
is positive. The output range is compressed towards the middle, always.

Three responses, all used in practice: rescale the output afterwards; use
\emph{mean of maxima} instead, which does reach the extremes but is discontinuous;
or accept it, on the grounds that a risk score is compared against other risk scores
rather than read as an absolute. Whichever you choose, know that the numbers a
Mamdani controller emits are not on the scale you think they are, and do not label
the output ``percent''.
\end{alertbox}

\section{Why Smoothness Is the Point}

Everything above could have been done with a threshold rule: \textsc{if} effort
$> 70$ \textsc{and} uncertainty $> 7$ \textsc{then} risk is high. It would be
shorter, and it has one fatal property.

\dsfig{%
\begin{tikzpicture}
\begin{axis}[
  width=12.2cm, height=5.6cm,
  xlabel={requirement uncertainty, with effort held at 45},
  ylabel={risk score},
  xlabel style={font=\scriptsize}, ylabel style={font=\scriptsize},
  tick label style={font=\tiny},
  xmin=0, xmax=10, ymin=-5, ymax=105,
  grid=both, major grid style={gray!18}, minor grid style={gray!8},
  legend style={font=\tiny, at={(0.03,0.97)}, anchor=north west,
                draw=gray!40, fill=white},
  legend cell align=left]

\addplot[greenhl, line width=1.4pt, mark=*, mark size=1.4pt] coordinates {
  (0,16.79) (1,32.71) (2,41.19) (3,46.21) (4,49.04) (5,49.75)
  (6,50.71) (7,53.50) (8,58.46) (9,66.81) (10,73.86)};
\addlegendentry{Mamdani controller}

\addplot[accent, line width=1.4pt] coordinates {
  (0,10) (7,10) (7,90) (10,90)};
\addlegendentry{threshold rule: high if uncertainty $>7$}

\node[font=\tiny\itshape, text=accent, anchor=west, align=left]
  at (axis cs:2.1,86)
  {one unit of measurement error here\\ changes the answer by 80 points};
\draw[accent!70, line width=0.7pt, -{Stealth[length=1.6mm]}]
  (axis cs:5.0,80) -- (axis cs:6.85,62);
\end{axis}
\end{tikzpicture}}%
{The same knowledge as a fuzzy controller and as a threshold rule. The fuzzy output
moves continuously; the threshold output is flat, flat, and then an eighty-point
jump between two adjacent inputs.}{fuzzy-smooth}

The threshold rule says the same thing at uncertainty $1$ and at $6$, and something
completely different between $7$ and $8$. If the uncertainty score is a judgement
with a unit of noise in it --- and it is, because a person assigned it --- then two
assessors will disagree wildly about projects near the boundary and not at all about
the rest. That is not a property anybody wants in a risk register.

The fuzzy controller's output moves a little when the input moves a little,
everywhere. \emph{That is the engineering reason fuzzy control was adopted}: not
because vagueness is philosophically interesting, but because interpolating between
rules gives a smooth response surface from a handful of sentences that domain
experts can actually write.

\begin{worked}[Beacon, by hand]
Effort $45$ person-days, requirement uncertainty $4$.

\smallskip
\textbf{Step 1 --- fuzzify.} With triangles peaking at $0$, $50$, $100$ for effort:
\[
  \mu_{\mathit{low}}(45) = \frac{50-45}{50} = 0.1, \qquad
  \mu_{\mathit{medium}}(45) = \frac{45-0}{50} = 0.9, \qquad
  \mu_{\mathit{high}}(45) = 0
\]
and with triangles peaking at $0$, $5$, $10$ for uncertainty:
\[
  \mu_{\mathit{low}}(4) = \frac{5-4}{5} = 0.2, \qquad
  \mu_{\mathit{medium}}(4) = \frac{4-0}{5} = 0.8, \qquad
  \mu_{\mathit{high}}(4) = 0
\]
Each input is described by exactly two words, as the partition guarantees.

\smallskip
\textbf{Step 2 --- rule strengths, by $\min$.} Five of the nine rules involve a word
with membership zero and contribute nothing. The other four:
\begin{tight}
  \item low effort, low uncertainty $\to$ \emph{low} risk: $\min(0.1, 0.2) = 0.1$
  \item low effort, medium uncertainty $\to$ \emph{low} risk: $\min(0.1, 0.8) = 0.1$
  \item medium effort, low uncertainty $\to$ \emph{low} risk: $\min(0.9, 0.2) = 0.2$
  \item medium effort, medium uncertainty $\to$ \emph{medium} risk:
        $\min(0.9, 0.8) = 0.8$
\end{tight}

\smallskip
\textbf{Step 3 --- clip.} Three rules conclude \emph{low} risk, with strengths
$0.1$, $0.1$ and $0.2$; aggregating by $\max$ means the \emph{low} triangle is
clipped at $0.2$. One rule concludes \emph{medium} risk, so that triangle is clipped
at $0.8$.

\smallskip
\textbf{Step 4 --- aggregate.} The combined shape is the pointwise maximum: a
plateau at height $0.2$ over the low-risk region and a larger plateau at $0.8$ over
the middle.

\smallskip
\textbf{Step 5 --- defuzzify.} The centre of area of that shape is
$\mathbf{49.04}$.

\smallskip
\textbf{Sanity check.} The dominant rule concluded \emph{medium}, whose triangle
peaks at $50$; the weaker \emph{low} contribution drags the answer slightly below
it. $49.04$ is exactly where it should be, and the arithmetic is worth doing once by
hand because it is the only way to believe the pipeline.

\smallskip
\textbf{Now move one input.} At uncertainty $9$ instead of $4$ the answer is
$66.81$; at uncertainty $0$ it is $16.79$. Every intermediate value appears ---
no jumps, and no rule had to be written for the cases in between.
\end{worked}

\begin{pitfall}
\begin{tight}
  \item \textbf{Reading a membership as a probability.} $\mu = 0.4$ is not ``a
        $40\%$ chance''. Memberships do not sum to one, do not update on evidence,
        and combine by $\min$ and $\max$ rather than by Bayes' rule.
  \item \textbf{Membership functions that do not overlap.} Then only one rule ever
        fires, there is no interpolation, and you have written a slow lookup table.
        Build a partition whose memberships sum to one.
  \item \textbf{Believing the output range.} Centroid defuzzification cannot reach
        the extremes --- here the output never leaves roughly $[16.6, 83.4]$. Never label
        it ``percent''.
  \item \textbf{Adding inputs.} Three linguistic variables with three values each
        need $27$ rules, four need $81$. Fuzzy controllers stay small for a reason.
  \item \textbf{Expecting entailment.} There is no soundness theorem here and
        $A \cap \neg A \ne \varnothing$. This is a control technique, not a logic in
        \chref{propositional}'s sense.
  \item \textbf{Using it where probability belongs.} If the question is ``will this
        release fail?'' the answer is uncertain, not vague, and \chref{uncertainty}
        is the right chapter.
\end{tight}
\end{pitfall}

%---------------------------------------------------------------------
\begin{lab}[Risk Rating by Mamdani Inference]
The full pipeline on five projects, the hand-checkable trace for \emph{Beacon}, the
smoothness comparison of \dsref{fuzzy-smooth}, and the output-range limitation.
\begin{lstlisting}[language=Python]
"""Chapter 16 lab -- rating project risk by Mamdani inference.

Two inputs, nine rules, one number out. The interesting outputs are
the smoothness table and the range check: a centroid defuzzifier
cannot reach the ends of its own scale.
"""

import numpy as np


def tri(x, a, b, c):
    """Triangular membership: 0 at a, 1 at b, 0 at c."""
    x = np.asarray(x, dtype=float)
    return np.clip(np.minimum((x - a) / (b - a), (c - x) / (c - b)), 0, 1)


# Fuzzy partitions: peaks evenly spaced, memberships summing to one,
# so every input is described by at most two words.
EFFORT = {                                   # person-days, 0..100
    "low":    lambda x: tri(x, -50, 0, 50),
    "medium": lambda x: tri(x, 0, 50, 100),
    "high":   lambda x: tri(x, 50, 100, 150),
}
UNCERTAINTY = {                              # the team's 0..10 scale
    "low":    lambda x: tri(x, -5, 0, 5),
    "medium": lambda x: tri(x, 0, 5, 10),
    "high":   lambda x: tri(x, 5, 10, 15),
}
RISK = {                                     # score, 0..100
    "low":    lambda x: tri(x, -50, 0, 50),
    "medium": lambda x: tri(x, 0, 50, 100),
    "high":   lambda x: tri(x, 50, 100, 150),
}

# (effort word, uncertainty word) -> risk word
RULES = [
    ("low", "low", "low"), ("low", "medium", "low"),
    ("low", "high", "medium"),
    ("medium", "low", "low"), ("medium", "medium", "medium"),
    ("medium", "high", "high"),
    ("high", "low", "medium"), ("high", "medium", "high"),
    ("high", "high", "high"),
]

GRID = np.linspace(0, 100, 1001)


def mamdani(effort, uncertainty, trace=False):
    # 1. fuzzify
    me = {k: float(f(effort)) for k, f in EFFORT.items()}
    mu = {k: float(f(uncertainty)) for k, f in UNCERTAINTY.items()}
    if trace:
        print("   effort %g      -> %s" % (effort,
              {k: round(v, 3) for k, v in me.items() if v > 0}))
        print("   uncertainty %g -> %s" % (uncertainty,
              {k: round(v, 3) for k, v in mu.items() if v > 0}))

    # 2. firing strength = min of the conditions
    # 3. clip the consequent at that strength
    # 4. aggregate by max
    aggregate = np.zeros_like(GRID)
    fired = []
    for (e, u, r) in RULES:
        w = min(me[e], mu[u])
        if w > 1e-9:
            fired.append((e, u, r, w))
            aggregate = np.maximum(aggregate, np.minimum(w, RISK[r](GRID)))
    if trace:
        for (e, u, r, w) in fired:
            print("   %-6s x %-6s -> %-6s   strength %.3f"
                  % (e, u, r, w))

    # 5. defuzzify by centroid
    area = aggregate.sum()
    crisp = float((GRID * aggregate).sum() / area) if area else 0.0
    return crisp, fired


def threshold_rule(effort, uncertainty):
    """The same knowledge as a crisp rule, for contrast."""
    return 90 if (effort > 40 and uncertainty > 7) else 10


PROJECTS = {
    "Atlas":  (20, 2),
    "Beacon": (45, 4),
    "Cobalt": (85, 9),
    "Delta":  (45, 9),
    "Echo":   (90, 2),
}

if __name__ == "__main__":
    print("project   effort  uncert    risk   rules fired")
    print("-" * 50)
    for name, (e, u) in PROJECTS.items():
        risk, fired = mamdani(e, u)
        print("%-9s %6d %7d %7.2f %8d" % (name, e, u, risk, len(fired)))
        assert len(fired) == 4, "a partition should fire 2x2 rules"

    print()
    print("=== Beacon in full (effort 45, uncertainty 4) ===")
    risk, fired = mamdani(45, 4, trace=True)
    print("   crisp risk = %.2f" % risk)
    assert abs(risk - 49.04) < 0.01

    # the memberships are exactly the hand calculation
    assert abs(float(EFFORT["low"](45)) - 0.1) < 1e-9
    assert abs(float(EFFORT["medium"](45)) - 0.9) < 1e-9
    assert abs(float(UNCERTAINTY["low"](4)) - 0.2) < 1e-9
    assert abs(float(UNCERTAINTY["medium"](4)) - 0.8) < 1e-9

    # --- smoothness ------------------------------------------------
    print()
    print("=== effort fixed at 45, uncertainty swept ===")
    print("   unc   fuzzy    threshold rule")
    print("   " + "-" * 34)
    curve = []
    for u in range(11):
        r, _ = mamdani(45, u)
        curve.append(r)
        print("   %3d %8.2f %12d" % (u, r, threshold_rule(45, u)))

    # the fuzzy output is monotone and moves in small steps
    steps = [curve[i + 1] - curve[i] for i in range(len(curve) - 1)]
    assert all(s > 0 for s in steps), "should be monotone increasing"
    print()
    print("   fuzzy: largest single step %.2f points" % max(steps))
    print("   threshold rule: one step of 80 points, at u = 7 -> 8")
    assert max(steps) < 20

    # --- the range limitation --------------------------------------
    print()
    print("=== what range can the output actually reach? ===")
    lo, hi = 1e9, -1e9
    for e in range(0, 101, 5):
        for u in range(0, 11):
            r, _ = mamdani(e, u)
            lo, hi = min(lo, r), max(hi, r)
    print("   across the whole input space, risk lies in [%.2f, %.2f]"
          % (lo, hi))
    print("   a minimum-effort, minimum-uncertainty project scores")
    print("   %.2f, not 0. Centroid defuzzification cannot reach the"
          % mamdani(0, 0)[0])
    print("   ends of its own scale, so never label this 'percent'.")
    assert lo > 5 and hi < 95

    # --- and the confusion this chapter exists to prevent ----------
    print()
    print("=== membership is not probability ===")
    print("   mu_high(effort=70) = %.3f" % float(EFFORT["high"](70)))
    print("   This does NOT mean a 40% chance the project is high")
    print("   effort. It is 70 person-days -- we are certain of that.")
    print("   The number describes the word 'high', not our ignorance,")
    print("   and no further measurement will change it.")
    total = sum(float(f(70)) for f in EFFORT.values())
    print("   memberships at 70 sum to %.3f (probabilities would sum"
          % total)
    print("   to 1 only by coincidence; here it is by construction)")
\end{lstlisting}
Then widen the \emph{medium} risk triangle to \texttt{tri(x, -25, 50, 125)} and
re-run. Beacon's score barely moves, but the reachable range narrows further: the
defuzzified output depends on the \emph{shape} of the output sets as much as on the
rules, which is the part of a fuzzy controller that is tuned rather than derived.
\end{lab}

%---------------------------------------------------------------------
\begin{checkpoint}
\begin{tightnum}
  \item A project is measured at exactly $70$ person-days and
        $\mu_{\mathit{high}}(70) = 0.4$. Explain why gathering more information
        about the project will not change that number, and what would.
  \item Four of nine rules fired for \emph{Beacon}, and the same four would fire for
        any input in a neighbourhood of it. Explain what property of the membership
        functions guarantees ``at most four'', and what goes wrong if the
        membership functions do not overlap.
  \item The controller's output never leaves roughly $[16.6, 83.4]$. Give the reason,
        and one consequence for how the score should be reported to a project
        board.
\end{tightnum}
\end{checkpoint}

%---------------------------------------------------------------------
\begin{chaptersummary}
\begin{tight}
  \item \term{Vagueness} is not \term{uncertainty}. A probability describes
        ignorance about a sharp fact and resolves with evidence; a \term{degree of
        membership} describes a blurry word and does not.
  \item A \term{fuzzy set} is a membership function $\mu_{A}(x) \in [0,1]$. A
        \term{fuzzy partition} --- overlapping triangles whose memberships sum to
        one --- means each input is described by at most two words and the
        description moves continuously.
  \item The operators are $\min$, $\max$ and $1-\mu$. Note that
        $A \cap \neg A \neq \varnothing$: there is no entailment relation and no
        soundness theorem, and this is a control technique rather than a logic.
  \item \term{Mamdani inference}: fuzzify, take each rule's strength as the $\min$
        of its conditions, clip its consequent at that strength, aggregate by
        $\max$, and defuzzify by \term{centroid}. For \emph{Beacon} four of nine
        rules fired and the answer was $49.04$.
  \item Centroid defuzzification \emph{cannot reach the extremes} --- the output
        here lives in roughly $[16.6, 83.4]$ whatever the input. Do not call it a
        percentage.
  \item The engineering justification is \emph{smoothness}: the output moves a
        little when the input moves a little, everywhere, where a threshold rule
        jumps eighty points between two adjacent inputs. A handful of sentences
        from a domain expert becomes a continuous response surface.
  \item Rules grow as $v^{n}$ in the number of inputs --- nine here, $81$ for four
        inputs --- which is why fuzzy controllers stay small.
\end{tight}
\end{chaptersummary}

\begin{reviewq}
  \item \textbf{(MCQ)} A degree of membership $\mu_{A}(x) = 0.4$ means:
        (a)~a $40\%$ chance $x$ is in $A$ (b)~$x$ belongs to $A$ to degree $0.4$
        (c)~$40\%$ of objects like $x$ are in $A$ (d)~$A$ is $40\%$ certain.
  \item \textbf{(MCQ)} In Mamdani inference a rule's firing strength is:
        (a)~the product of its condition memberships (b)~the $\min$ of them
        (c)~the $\max$ of them (d)~their average.
  \item \textbf{(MCQ)} Aggregation across rules combines the clipped consequents
        by: (a)~$\min$ (b)~$\max$ (c)~addition (d)~multiplication.
  \item \textbf{(MCQ)} Centroid defuzzification: (a)~always reaches the extremes of
        the output range (b)~compresses the output towards the middle
        (c)~is discontinuous (d)~requires probabilities.
  \item \textbf{(Short)} Give one question about a software release that is a
        matter of uncertainty and one that is a matter of vagueness, and say which
        chapter handles each.
  \item \textbf{(Short)} Explain why membership functions must overlap, and what a
        fuzzy controller degenerates into if they do not.
  \item \textbf{(Short)} State the engineering argument for a fuzzy controller over
        a threshold rule, using \dsref{fuzzy-smooth}.
  \item \textbf{(Applied)} Compute by hand the fuzzy output for a project with
        effort $75$ and uncertainty $6$: give the memberships, the firing strengths
        of every rule that fires, which output words are clipped and at what level,
        and say whether the crisp result should be above or below $50$, with a
        reason.
\end{reviewq}
